\chapter{Tipos y operadores en C-Horizon--}
\label{cap:tipos y operadores}
En esta sección se especifican los tipos y operadores que vamos a usar. Los operadores infijos serán señalados como tal. En caso de empate de prioridades,
siempre asociaremos por la izuierda.

\section{Integer}
Corresponde con el \textit{signed int} de C++. Es decir, almacena números enteros representados en sistema decimal. Tienen un espacio reservado de 32 bits, 
por lo que su rango de valores es [-2147483648, 2147483647]. La palabra reservada para este tipo es  \Integer.
 
\section{Boolean}
Corresponde con el \textit{bool} de C++. Almacenará valores verdaderos (\const{True}) o falsos (\const{False}). Se realizará conversión automática de tipo 
\Integer{} siendo "$0$" el entero que represente falso y cualquier otro número, verdadero. Las operaciones que podremos realizar serán:
 

%OPERADORES ARITMÉTICOS
\section{Operadores aritméticos (para enteros)}

\subsection*{Suma}
Suma entera de dos números enteros. Se utilizará el operador infijo \Sum. \\Tipado: \Sum :: (\Integer, \Integer) $\rightarrow$ \Integer. \\Prioridad: 7.

\subsection*{Resta}
Resta entera de dos números enteros. Se utilizará el operador infijo \Sub. \\Tipado: \Sub :: (\Integer, \Integer) $\rightarrow$  \Integer. \\Prioridad: 7.

\subsection*{Multiplicación}
Multiplicación entera de dos números enteros. Se utilizará el operador \Mul. \\Tipado: \Mul :: (\Integer, \Integer) $\rightarrow$  \Integer. \\Prioridad: 8.

\subsection*{División}
División entera de dos números enteros. Se truncará el resultado obtenido en caso de ser fraccionario. El operador utilizado será \Div. \\Tipado: \Div :: 
(\Integer, \Integer) $\rightarrow$  \Integer. \\Prioridad: 8.

%OPERADORES LÓGICOS
\section{Operadores lógicos (para booleanos)}

\subsection*{Conjunción}
Conjunción lógica de dos booleanos. Su operador infijo será el \And. \\Tipado: \And :: (\Boolean, \Boolean) $\rightarrow$ \Boolean. \\Prioridad: 4.

\subsection*{Disyunción}
Disyunción lógica de dos booleanos. Su operador infijo será el \Or. \\Tipado: \Or :: (\Boolean, \Boolean) $\rightarrow$ \Boolean. \\Prioridad: 5.

\subsection*{Negación}
Negación lógica de dos booleanos. Su operador infijo será el \Not. \\Tipado: \Not :: \Boolean $\rightarrow$ \Boolean. \\Prioridad: 6.

%OTROS OPERADORES
\section{Operadores genéricos}

\subsection*{Igual que}
Compara dos valores de un mismo tipo y devuelve \const{True} si tienen el mismo valor, en caso contrario devuelve \const{False}. \\Tipado: \Eq :: (t, t)
\footnote{Usaremos "t" para designar un tipo cualquiera} $\rightarrow$ \Boolean. \\Prioridad 3.

\subsection*{Distinto que}
Equivale a \Not (a \Eq b). \\Prioridad 3.

\subsection*{Menor que}
Compara dos valores de un mismo tipo y devuelve \const{True} si el primero es menor que el segundo, en caso contrario devuelve \const{False}
\footnote{Consideramos que \const{False}  es menor que \const{True}}. Su operador infijo es \Lt. 
\\Tipado: \Lt :: (t, t) $\rightarrow$ \Boolean. \\Prioridad: 3.

\subsection*{Mayor que}
Compara dos valores y devuelve \const{True} si el primero es mayor que el segundo, en caso contrario devuelve \const{False}. Su operador infijo es \Gt.
\\Tipado: \Gt :: (t, t) $\rightarrow$ \Boolean. \\Prioridad: 3.

\subsection*{Menor o igual que}
Equivale a \Not (a \Lt b). \\Prioridad: 3.

\subsection*{Mayor o igual que}
Equivale a \Not (a \Gt b). \\Prioridad: 3.

\subsection*{Asignación}
Asigna un valor a una variable. El tipo del valor asignado debe corresponder con el tipo de la variable. Su operador infijo será ":="
(sin las dobles comillas). \\Tipado: ??.\\Prioridad: 1.
\begin{algorithm}
	\Integer{} num := 5; \comm{//Asignación válida} \\
	\Integer{} num2 := \const{True}; \comm{//Asignación no válida por fallo de tipos} \\
\end{algorithm}

\subsection*{Struct}
Sirve para declarar registros.

\subsection*{Class}

\subsection*{Paso por valor}
Indica que el argumento de la función se pasa por valor. Lo que realiza realmente es crear una nueva variable con el valor de la anterior. Esta variable se crea
en el ámbito de la función, de tal manera que se destruye al terminar ésta. Su operador será \Val.
\\Tipado: \Val :: t $\rightarrow$ t. \\Prioridad 1.

\subsection*{Punteros}
Indica que el argumento de la función se pasa por valor. Lo que realiza realmente es crear una nueva variable con el valor de la anterior. Esta variable se crea
en el ámbito de la función, de tal manera que se destruye al terminar ésta. Su operador será \Val.
\\Tipado: \Val :: t $\rightarrow$ t. \\Prioridad 1.


%TABLA RESUMEN DE PRIORIDADES
